\begin{figure}[!b]
\centering
\includestandalone[mode=buildmissing]{Standalone/Plots/TPUvsPLA_neg_deg}
    \caption{Resistance in varying environmental conditions of two sensors on a TPU substrate, one with a conductive-PLA trace and one with a conductive-TPU trace: a) resistance and temperature versus time, b) PLA resistance against ambient temperature showing the hysteresis, c) TPU resistance against ambient temperature showing the hysteresis.}
    \label{fig:TPUvsPLA_neg_deg}
\end{figure}

The resistance of both conductor materials responds strongly to temperature, which is critical for the present application since the sensor must operate reliably across a range of \SIrange{-20}{50}{\celsius}. To characterize this behavior, two sensors on a TPU substrate, one with a conductive-PLA trace and one with a conductive-TPU trace, were placed in a climate chamber and subjected to repeated temperature cycles while their resistance was recorded. The results are shown in \cref{fig:TPUvsPLA_neg_deg}, where panel a) gives the full resistance and temperature trace versus time, and panels b) and c) show the resistance of each sensor plotted against the ambient temperature to reveal the hysteresis in the response.

Both sensors exhibit a strong positive dependence on temperature, with the resistance rising as the temperature increases and falling as it decreases. For the conductive-PLA sensor, the resistance rises from approximately \SI{55}{\kilo\ohm} at \SI{-20}{\celsius} to approximately \SI{115}{\kilo\ohm} at \SI{50}{\celsius}, a factor of roughly two across the operating range. The conductive-TPU sensor spans a similar factor, from approximately \SI{90}{\kilo\ohm} to approximately \SI{165}{\kilo\ohm} over the same temperature range. This temperature dependence is large relative to the strain signal and would dominate the output of a single sensor if left uncompensated, which motivates the Wheatstone bridge configuration discussed in \Cref{Wheatstone Half-Bridge Design}. 
The hysteresis plots in panels \cref{fig:TPUvsPLA_neg_deg}b) and c) reveal a further difference between the two materials. The conductive-PLA sensor, shown in panel b), traces a relatively narrow loop when the temperature cycles, indicating that the resistance at a given temperature differs only moderately between the heating and cooling directions. The conductive-TPU sensor, in panel c), exhibits considerably wider loops. These wide loops are consistent with the pronounced drift observed for the conductive TPU in \Cref{Material & Mounting-Surface Selection (Drift)}.

The concern raised that conductive PLA might not tolerate temperatures at the extremes of, or beyond, the intended operating range was investigated by exposing the same PLA-based sensors to \SI{-20}{\celsius}, \SI{-40}{\celsius}, \SI{60}{\celsius}, \SI{70}{\celsius}, and \SI{80}{\celsius}, after which their strain response was measured at room temperature, with the results shown in \cref{fig:PLA_after_temp_limits}. This approach, in which the sensor is first exposed to a given temperature and only afterwards strained, was necessary because the available test setup could not apply strain inside the climate chamber across the extreme temperatures of interest. The strain response is largely preserved following all temperature exposures, where both the resistance and $\Delta R / R_0$ curves closely follow the reference measurement, the baseline resistance shifts only slightly for some exposures, and only a small spread appears at the peak response. These results indicate that exposure to extreme temperatures does not measurably affect the sensitivity of an unstrained sensor.

\begin{figure}[!b]
\centering
\includestandalone[mode=buildmissing]{Standalone/Plots/PLA_efter_neg_deg}
    \caption{Resistance of a PLA sensor exposed to strain after different temperature exposures. The upper panel shows the measured resistance and the lower panel the relative change $\Delta R/R_0$.}
    \label{fig:PLA_after_temp_limits}
\end{figure}

The test does not demonstrate that the sensors can withstand these temperatures while simultaneously being strained. The strains encountered in the intended application are, however, very small, so the untested condition represents only a mild mechanical demand. Combined with the preserved strain response observed after temperature exposure, this provides evidence that the sensors are likely to tolerate operation across this temperature range without degradation, although this cannot be conclusively verified from the present measurements.
Taking the temperature results together with the drift measurements presented in \Cref{Material & Mounting-Surface Selection (Drift)}, conductive PLA is the preferred material for this application. Its lower drift, narrower temperature-induced hysteresis, and preserved strain response following exposure to extreme temperatures all represent clear advantages over conductive TPU. The sensors exhibited a stable and repeatable response both during temperature cycling and after low and high temperature exposure, with no indication of mechanical failure or loss of sensitivity. Conductive PLA was therefore selected as the sensing material for all subsequent measurements.